\subsection*{Основные операции:}

\textbf{Insert: }
 \begin{lstlisting}
void Heap::Insert(int x) {
  tree.push_back(x); 
  ++size; 
  SiftUp(size - 1); 
}
\end{lstlisting}
Вставляем элемент в конец. Тогда могли нарушиться только неравенства с родителем. Поэтому необходимо просеять наверх.
\textit{Асимптотика: O(log(n)) за счет выполнения операции SiftUp()}

\textbf{GetMin: }
 \begin{lstlisting}
int Heap::GetMin() { 
  return tree[0]; 
}
\end{lstlisting}
Минимальный элемент кучи лежит в вершине дерева, просто его выводим. 
\textit{Асимптотика: O(1)}

\textbf{ExtractMin: }
 \begin{lstlisting}
void Heap::ExtractMin() {
  tree[0] = tree[size - 1]; 
  tree.pop_back();
  --size;
  SiftDown(0);
}
\end{lstlisting}
Меняем первый элемент и последний местами, удаляем последний. Т.к. у первого элемента нет родителей, то неравенства могли нарушиться только с детьми, поэтому просеиваем элемент вниз.\\
\textit{ Асимптотика: O(log(n)) за счет выполнения операции SiftDown()}

\pagebreak